\documentclass[10pt]{article}
\usepackage[margin=20mm]{geometry}
\usepackage{amsmath,amssymb,mathtools,fontspec,unicode-math,xurl,hyperref}
\IfFontExistsTF{DejaVu Serif}{\setmainfont{DejaVu Serif}}{\setmainfont{Latin Modern Roman}}
\setmathfont{latinmodern-math.otf}
\setlength{\parindent}{0pt}\setlength{\parskip}{6pt}\sloppy
\title{Technical Review of
The Contact Process on Dynamic Regular Graphs:
Subcritical Phase and Monotonicity}
\author{AI-assisted technical review}
\date{}
\begin{document}
\maketitle
Abstract: The reviewed version contains two independently demonstrated defects in the written strict-monotonicity argument: the passage from split edges to successor states loses multiplicities and, at positive parameter increment, changes the evolution of untouched herds; a subsequent reciprocal bound uses a minimum where a maximum is required. Exact local generator calculations corroborate the first finding. A separate discretization error omits a terminal time interval in the birth-count estimate. These are errors in displayed identities or proof steps, not counterexamples to the final theorems. Explicit repairs are given within the checked scope. Scientific adjudication remains pending.\par
\section*{Paper identity and scope}
The paper is by Bruno Schapira and Daniel Valesin, The Annals of Probability 53(2) (2025), 753–796, DOI 10.1214/24-AOP1721 (https://doi.org/10.1214/24-AOP1721). The reviewed file is the 44-page publisher-formatted version of record acquired from the author’s website. Printed page \(p\) is physical PDF page \(p-752\). All page locators below refer to printed pages unless explicitly stated otherwise.\par
The complete extracted text, including the appendix and references, was read. The most intensive independent checks concern Definitions 2.1–2.4, the particle-count generator, Proposition 3.7, Lemma 3.8, the use of those results in Proposition 3.6, and the time partition in Corollary 3.15. The source pages supporting these findings were also viewed as rendered images. Reading the other proofs is not represented as an independent certification of every argument or of the external results they cite.\par
The author publication list and arXiv version history were consulted for correction and version context on 29 October 2026. The inspected arXiv record 2309.17040 lists a single 65-page version dated 29 September 2023; it is not the reviewed 44-page journal file. The inspected author entry links the journal-formatted article. No relevant correction was located in this bounded check. This does not establish that no correction exists elsewhere. No author code was executed.\par
\section*{Definitions and operative assumptions}
Fix an integer \(d\geq3\), birth rate \(\lambda>0\), and splitting rate \(v>0\). A herd shape \(A\) is any finite nonempty subset of the infinite \(d\)-regular tree. It need not be connected (Definition 2.1, p. 756). A state \(\xi\in\mathcal S\) is a finite multiset of such shapes. For an edge \(e\), write \(A_{e,1},A_{e,2}\) for the portions on the two sides of \(e\). An edge is active precisely when both portions are nonempty.\par
Equation reference: gen. Definition 2.3, p. 757, gives a death rate of one per particle, birth rate \(\lambda\) per occupied-neighbor arrow to an unoccupied site of the same herd, and split rate \(v\) per active edge per herd copy. Different herds evolve independently, even when their spatial locations overlap. Let \(X(\xi)\) be the total particle count, \(H(\xi)\) the total number of occupied adjacent pairs within herds, and \(\mathcal E(\xi)\) the active-edge count with herd multiplicities. Let \(L_v\) denote this generator and put \[S_v(t)\xi=\mathbb E_{\lambda,v}[X_t\mid\Xi_0=\xi],
 \qquad F_v(t)=S_v(t)\delta_{\{o\}}.\] Splitting preserves \(X\). Counting the available birth arrows gives \[
 L_vX=(d\lambda-1)X-2\lambda H .\] The paper’s pure-birth domination and Corollary 2.3 provide finite moments on bounded time intervals. Proposition 2.7 implies \(F_v(t+s)\leq F_v(t)F_v(s)\). These facts supply the integrability and positivity used below; in particular, \(F_v(t)\geq e^{-t}>0\), by survival of the initial particle.\par
\section*{Errors in the derivative kernel}
Problem 1 (Split multiplicity and the finite-increment kernel).  The successor-state rewriting in (28)–(29), pp. 769–770, and the associated first-separation calculation in Lemma 3.10, p. 771, do not follow from the stated definitions. Severity: Major, as an invalid step in a core derivation; no false final theorem is established.\par
\subsection*{Original assertions}
Equation reference: g. For every \(\xi\in\mathcal S\) and \(t\geq0\), (20) defines \[
 g_v(\xi,t)=\sum_A\xi(A)\sum_{e\text{ active in }A}
 \bigl[S_v(t)(\delta_{A_{e,1}}+\delta_{A_{e,2}})
                  -S_v(t)\delta_A\bigr].\] Equation (27) replaces the first \(v\) in each bracket by \(v+\varepsilon\), where \(\varepsilon>0\), and calls the result \(g_{v,\varepsilon}\). Definition 3.2 defines \(G(\xi)\) as a set of successor configurations and, for the shape \(A\) split to obtain \(\xi'\), sets \(m(\xi,\xi')=\xi(A)\). Equations (28) and (29) assert, respectively, that the preceding edge sums equal \[\sum_{\xi'\in G(\xi)}m(\xi,\xi')
       [S_v(t)\xi'-S_v(t)\xi],\qquad
 \sum_{\xi'\in G(\xi)}m(\xi,\xi')
       [S_{v+\varepsilon}(t)\xi'-S_v(t)\xi].\] There are two distinct failures.\par
\subsection*{Two active edges can have the same successor}
Let \(a,b,c\) be three successive vertices of the tree and \(A=\{a,c\}\), \(\xi=\delta_A\). Both \(ab\) and \(bc\) are active. They give the same successor \(\xi'=\delta_{\{a\}}+\delta_{\{c\}}\). Thus \(G(\xi)\) has one element and the printed \(m\) is one, whereas Definition 2.3 gives the transition \(\xi\to\xi'\) total rate \(2v\). Likewise the extra split clock at rate increment \(\varepsilon\) gives first-order probability \(2\varepsilon s+o(s)\), rather than \(\varepsilon s+o(s)\).\par
This state satisfies all the definitions. It also occurs with positive probability starting from a singleton: require two consecutive births along \(a,b,c\), then death of \(b\), with no intervening split or other event. All prescribed transition rates are positive and the finite list of holding rates is finite. Restricting attention to the actual process started at \(\delta_{\{o\}}\) therefore does not remove the example.\par
Equation reference: d2. The discrepancy in (28) is substantive even as an identity of expectations. Set \[D(t)=S_v(t)\xi'-S_v(t)\xi.\] Both configurations have \(X=2,H=0\), so \(D(0)=D'(0)=0\). At \(\xi\), birth at the middle vertex has rate \(2\lambda\) and adds two occupied adjacent pairs; the other \(2(d-1)\) birth arrows each add one. At \(\xi'\), there are \(2d\) arrows, each adding one within its own herd. Consequently \[L_vH(\xi)=2(d+1)\lambda,\qquad L_vH(\xi')=2d\lambda .\] Applying [gen] (\#gen) twice yields \[
 D''(0)=4\lambda^2,\qquad
 D(t)=2\lambda^2t^2+o(t^2)>0
 \quad\text{for all sufficiently small }t>0.\] The edge definition is \(2D(t)\); the printed state sum is \(D(t)\).\par
The use of right derivatives at zero is justified by the finite local transition lists and the bounded-time moment domination: the finite generator iterations here have polynomial growth in the relevant counts, and the finite-step expansion can be localized and then dominated. An independent exact rational enumeration confirms [d2] (\#d2). For \(d=3,\lambda=1/6,v=1\), its results are \[\begin{array}{c|rrr}
 &X&L_vX&L_v^2X\\ \hline
 \delta_{\{a,c\}}&2&-1&1/18\\
 \delta_{\{a\}}+\delta_{\{c\}}&2&-1&1/6
\end{array}\] and hence \(D''(0)=1/9\), \(D(t)=t^2/18+o(t^2)\). The enumeration aggregates all elementary transition rates, checks full degree at every evaluated vertex, and uses rational arithmetic. It is a local generator calculation, not a simulation or a numerical approximation to a large-time theorem.\par
\subsection*{Untouched herds do not cancel at positive increment}
Equation reference: untouched. Even after multiplicities are fixed, (29) is incorrect with the definition (27). Take \(\xi=\delta_{\{a,b\}}+\delta_{\{z\}}\), where \(a\sim b\) and \(z\) is any vertex. There is exactly one active edge, so the preceding counting problem is absent. The only split replaces \(\{a,b\}\) by \(\{a\},\{b\}\). By additivity across herds, the right-hand side of (29) minus (27) is \[
 F_{v+\varepsilon}(t)-F_v(t).\] It is not zero in general. A direct generator calculation at a singleton gives \[\left.\frac{d^k}{dt^k}
   (F_{v+\varepsilon}(t)-F_v(t))\right|_{t=0}=0
 \quad (k=0,1,2),\qquad
 \left.\frac{d^3}{dt^3}
   (F_{v+\varepsilon}(t)-F_v(t))\right|_{t=0}
       =2d\lambda^2\varepsilon .\] To see the last coefficient, write \(c=d\lambda-1\) and \(L_v^2X=cL_vX-2\lambda L_vH\). A singleton gives an adjacent pair at total rate \(d\lambda\). At that pair, splitting decreases \(H\) by one, so the \(v\)-dependent part of \(L_v^2H\) at the singleton is \(-d\lambda v\). No lower-order term depends on \(v\). This proves the displayed difference.\par
The exact enumeration with \(d=3,\lambda=1/6,v=1,\varepsilon=1\) returns \[\begin{array}{c|rrrr}
 &F(0)&F'(0)&F''(0)&F'''(0)\\ \hline
 v=1&1&-1/2&1/12&43/72\\
 v=2&1&-1/2&1/12&55/72
 \end{array}\] and [untouched] (\#untouched) is \(t^3/36+o(t^3)>0\). Definition 3.3 evolves the whole second process with parameter \(v+\varepsilon\) after separation. Thus the omitted singleton term is also present in the actual conditional calculation used in Lemma 3.10.\par
\subsection*{Alternative readings, correction, and impact}
Interpreting \(G\) as a multiset indexed by elementary split events would address the first problem, but that is a change to the explicit set definition, not the printed identity. Assuming connected shapes contradicts Definition 2.1 and is not invariant under deaths. Neither reading addresses the untouched-herd term at \(\varepsilon>0\).\par
Equation reference: correctkernel. One may retain edge indexing throughout and use \[
 \widetilde g_{v,\varepsilon}(\xi,t)
 =\sum_A\xi(A)\sum_{e\text{ active in }A}
 \left[S_{v+\varepsilon}(t)
    \bigl(\xi-\delta_A+\delta_{A_{e,1}}+\delta_{A_{e,2}}\bigr)
       -S_v(t)\xi\right].\] Alternatively the correct successor multiplicity is the total number of elementary split events, including copies, that yield that successor. At \(\varepsilon=0\), additivity cancels untouched herds in [correctkernel] (\#correctkernel), giving exactly [g] (\#g). For \(0\leq t\leq T\), pure-birth domination gives the uniform bound \[|\widetilde g_{v,\varepsilon}(\xi,t)|
 \leq 2e^{d\lambda T}X(\xi)\mathcal E(\xi).\] This supplies the same kind of integrability needed in Lemmas 3.9–3.11 and the appendix, pp. 793–794. Finite-time parameter continuity gives \(\widetilde g_{v,\varepsilon}\to g_v\); the extra-split separation probability tends to zero as \(\varepsilon\downarrow0\), by the bounded-time active-edge moment bound. These observations support replacing the finite-increment kernel and repeating the dominated-convergence argument.\par
The error affects the written proof of Proposition 3.7, which supplies the derivative estimate in Proposition 3.6 and hence Theorem 1.3. The subcritical characterization used in Section 3.4 also invokes Proposition 3.6. This is why the finding concerns a core derivation. Equation (20) itself has the correct elementary-event indexing; the limiting derivative formula and the final theorems have not been disproved. The proposed correction is a mathematical repair identified here, not an author-adopted correction or an independently accepted revision.\par
\section*{The lower bound used for strict monotonicity}
Problem 2 (Reciprocal extremum and a time-range restriction).  The last display in the proof of Lemma 3.8, p. 768, is false under the paper’s assumptions. The subsequent use of this lemma in Proposition 3.6, p. 769, also exceeds its time range. Severity: Major for the stated core lower-bound derivation; both identified steps admit local repairs.\par
\subsection*{The false inequality and a counterexample}
Equation reference: wrongmin. Lemma 3.8 compares one adjacent two-particle herd with two singleton herds, for all \(t\geq1\), asserting a gap at least \(\gamma F_v(t)\), where \(\gamma>0\) depends continuously on \(\lambda,v\). The coupling event in its proof first gives a lower bound involving \(\min_{t-1\leq s\leq t}F_v(s)\). The last display asserts \[
 \min_{t-1\leq s\leq t}F_v(s)
 \geq F_v(t)\left(\min_{0\leq r\leq1}F_v(r)\right)^{-1}.\] Take the admissible parameters \(d=3,\lambda=1/6,v=1\). Equation [gen] (\#gen) and \(H\geq0\) imply \(F_v'(t)\leq-F_v(t)/2<0\). Consequently \(0<F_v(1)\leq e^{-1/2}<1\). At \(t=1\), the left side of [wrongmin] (\#wrongmin) is \(F_v(1)\), while the right side is \(F_v(1)/F_v(1)=1\). This is a strict contradiction. It uses the actual process, not an arbitrary submultiplicative test function.\par
\subsection*{Repair and its limits}
Submultiplicativity correctly implies \[F_v(s)\geq \frac{F_v(t)}{F_v(t-s)}
 \geq \frac{F_v(t)}{M_v},\qquad
 M_v=\max_{0\leq r\leq1}F_v(r),\quad t-1\leq s\leq t .\] Thus the minimum in the denominator must be replaced by a maximum. If \(q>0\) denotes the probability of the coupling event actually used in the proof, \(\gamma=q/M_v\) is a valid choice. Here \(M_v\) is finite and positive. Joint finite-time continuity, obtained by the same bounded-time coupling and moment domination, makes \(M_v\) continuous on the parameter domain. The prescribed finite exponential clock event has positive continuous probability. The stated positivity and continuous dependence can therefore be retained.\par
The proof of Proposition 3.6 applies Lemma 3.8 with time argument \(t-s\) but integrates over \(1\leq s\leq t\). The lemma only permits \(t-s\geq1\). At \(s=t\), the split and unsplit particle counts agree, so a positive multiple of \(F_v(0)=1\) cannot be a valid gap. Restrict the integration to \(1\leq s\leq t-1\). For \(t\geq3\), its length \(t-2\) is at least \(t/3\); submultiplicativity gives \[F_v(s-1)F_v(t-s)\geq F_v(t-1)
       \geq F_v(t)/F_v(1).\] Accordingly the needed lower bound by a positive constant times \(tF_v(t)\) survives for all sufficiently large \(t\), with a changed constant. This is sufficient for the growth-index limit used there. Together with a corrected derivative proof, these checks preserve the specific strict-monotonicity strategy. They do not turn the erroneous display into a true one, nor certify all later arguments.\par
\section*{A terminal interval in the birth-count estimate}
Problem 3 (Floor versus ceiling in Corollary 3.15).  The union bound on p. 776 omits births in the final fractional time interval. Severity: Minor; the stated moment conclusion has a direct repair.\par
Under the section’s assumptions \(\lambda<\bar\lambda(v)\) and \(0<\varphi<1\), the proof fixes \(p\geq1\), \(x>1\), and \(t_x=(p+1)\log x/|\log\varphi|\). It bounds \(\mathbb P(N_{t_x}>x)\) by a sum over \(k=0,\ldots,\lfloor t_x\rfloor-1\) of events involving \(N_{k+1}-N_k\). For a concrete admissible choice of the time, set \(x=\exp(|\log\varphi|/[2(p+1)])\), so \(t_x=1/2\). The displayed sum is empty. Its left side is positive: the finite prescribed event of more than \(x\) successive births along a tree path before time \(1/2\), with no interfering death or split, has positive probability. All rates required for that event are positive and finite. This demonstrates a false bound, including within the stated \(x>1\) range.\par
Let \(m=\lceil t_x\rceil\). Monotonicity of \(N\) and the union bound instead give \[\mathbb P(N_{t_x}>x)
 \leq \sum_{k=0}^{m-1}\mathbb P(N_{k+1}-N_k>x/m)
 \leq m^{p+2}x^{-(p+1)}
       \sup_{t\geq0}\mathbb E[(N_{t+1}-N_t)^{p+1}].\] The supremum is the one controlled by (37) and (44). For \(x\geq e\), \(m=O(\log x)\), so the tail estimate needed to integrate the \(p\)th moment is unchanged in order. The bounded range of \(x\) contributes a finite amount. This repair establishes the specific partition step without altering the claim that all finite positive moments of the total birth count exist in the subcritical regime. It is not evidence against that conclusion.\par
\section*{Assessment and coverage limits}
The most consequential finding is the mismatch between the finite-increment coupling and its kernel. The counterexamples satisfy the stated domains, retain the exact positive parameters, and cannot be removed by an implicit connectivity convention. The reciprocal and partition errors are separately demonstrated and have explicit corrections. The identified repairs give substantial counterevidence against interpreting these defects as failures of the final theorems.\par
The review also checked the direction of the particle-count domination, the role of herd independence in expectation additivity, the submultiplicative inequality, and the dependence of the growth-index argument on the derivative formula. The finite-graph coupling in Section 4, all higher-moment constants, and the results imported from earlier papers were read but were not independently re-proved in full. No global correctness certificate is claimed.\par
Dimension | Recorded assessment\par
Core proof defects established | Problems 1 and 2 (Major)\par
Local discretization defect | Problem 3 (Minor)\par
Final theorem counterexample | Not established\par
Author adoption of proposed repairs | Not established\par
Scientific adjudication | Pending\par
Independent human reviews | None recorded\par
\section*{Supporting checks}
The reviewer-authored calculation `generator\_check.py` enumerates deaths, births, and every active-edge split on a radius-six ball in the 3-regular tree. Only generator powers through order three at the specified local states are evaluated. Every transition evaluation asserts full degree at all occupied vertices, so no boundary transition enters these results. Exact fractions are used throughout. The saved output verifies the two derivative tables and both distinct-edge multiplicities. This code was written and executed for this review; it is not an author program.\par
`789ef90a43ed3a97a1f1d88d6e9312f7cfc7a1a90d08ed8ce1ca0d58ec678632`.\par
The scientific reading, derivations, report, and image inspection were performed in the current AI-assisted review conversation. A bounded native helper checked file integrity and compiler capabilities only. Its mechanical work is not scientific review or independent acceptance.\par
9 B. Schapira and D. Valesin, The contact process on dynamic regular graphs: Subcritical phase and monotonicity, Ann. Probab. 53 (2025), 753–796. doi:10.1214/24-AOP1721 (https://doi.org/10.1214/24-AOP1721). Reviewed author-hosted version of record: https://math.univ-lyon1.fr/\textasciitilde{}schapira/articlespdf/AOP.herdcontact.pdf. B. Schapira and D. Valesin, arXiv:2309.17040, submission history, https://arxiv.org/abs/2309.17040. Metadata consulted 29 October 2026; the 65-page preprint is not substituted for the reviewed journal file. B. Schapira, publication list, https://math.univ-lyon1.fr/\textasciitilde{}schapira/recherche.html. Consulted 29 October 2026 for the paper entry and bounded correction search.\par
\end{document}
